\begin{enumerate}
\item On considère les suites $\left(a_n\right)$ et $\left(b_n\right)$ définie par $a_0 = 1$ et, pour tout entier naturel $n$,\: $a_{n+1} = 0,5a_n + 1$ et $b_n = a_n - 2$.

On peut affirmer que:

\begin{center}
\begin{tabularx}{\linewidth}{*{2}{X}}
\textbf{a.~} $\left(a_n\right)$ est arithmétique ;&\textbf{b.~} $\left(b_n\right)$ est géométrique;\\
\textbf{c.~} $\left(a_n\right)$ est géométrique;&\textbf{d.~} $\left(b_n\right)$ est arithmétique. 
\end{tabularx}
\end{center}

\begin{tabular}{@{\hspace{0.05\linewidth}} | p{0.9\linewidth}}
$b_n=a_n-2$ donc $a_n=b_n+2$.

$b_{n+1}=a_{n+1}-2 = 0,5 a_n +1-2 = 0,5\left ( b_n+2\right ) -1
=0,5 b_n +1-1=0,5 b_n$ \\
donc la suite $(b_n)$ est géométrique de raison $0,5$.
\hfill\textbf{Réponse b.}
\end{tabular}
\end{enumerate}

\medskip

Dans les questions 2. et 3., on considère les suites $\left(u_n\right)$ et $\left(v_n\right)$ définies par :

\[u_0 = 2,\: v_0 = 1 \:\: \text{et, pour tout entier naturel }\:\:n : 
\left\{\begin{array}{l !{=} l}
u_{n+1}&u_n + 3v_n\\
v_{n+1}&u_n + v_n.
\end{array}\right.\]

\begin{enumerate}[resume]
\item On peut affirmer que :

\begin{center}
\begin{tabularx}{\linewidth}{*{4}{X}}
\textbf{a.~}$\left\{\begin{array}{l !{=} l}
u_2&  5\\ v_2 & 3
\end{array}\right.$&\textbf{b.~} $u_2^2 - 3v_2^2 = - 2^2$&\textbf{c.~} $\dfrac{u_2}{v_2} = 1,75$&\textbf{d.~} $5u_1 = 3v_1$.\\
\end{tabularx}
\end{center}

\begin{tabular}{@{\hspace{0.05\linewidth}} | p{0.9\linewidth}}
$\left\{\begin{array}{l !{=} l  !{=} l  !{=} l}
u_1 &  u_0+ 3v_0 & 2+3\times 1 & 5\\ 
v_1 & u_0+v_0 & 2 +1 & 3
\end{array}\right.$
\hfill
$\left\{\begin{array}{l !{=} l  !{=} l  !{=} l}
u_2 &  u_1+ 3v_1 & 5+3\times 3 & 14\\ 
v_2 & u_1+v_1 & 5 +3 & 8
\end{array}\right.$

$\dfrac{u_2}{v_2}=\dfrac{14}{8}=1,75$
\hfill\textbf{Réponse c.}
\end{tabular}

\medskip

\item On considère le programme ci-dessous écrit en langage Python :

\begin{center}
\begin{tabular}{|>{\texttt}l|}\hline
\texttt{def valeurs() :}\\
\quad \texttt{u = 2}\\
\quad \texttt{v = 1}\\
\quad \texttt{for k in range(1,11) :}\\
\quad \quad \texttt{c = u}\\
\quad \quad \texttt{u = u + 3*v}\\
\quad \quad \texttt{v = c + v}\\
\quad \texttt{return (u,v)}\\ \hline
\end{tabular}
\end{center}

Ce programme renvoie :

\begin{center}
\begin{tabularx}{\linewidth}{*{4}{X}}
\textbf{a.~} $u_{11}$ et $v_{11}$ ;&
\textbf{b.~} $u_{10}$ et $v_{11}$ ;&
\textbf{c.~} les valeurs de $u_n$ et $v_n$ pour $n$ allant de 1 à 10;&
\textbf{d.~} $u_{10}$ et $v_{10}$.
\end{tabularx}
\end{center}

\begin{tabular}{@{\hspace{0.05\linewidth}} | p{0.9\linewidth}}
\hfill\textbf{Réponse d.}
\end{tabular}
\end{enumerate}

Pour les questions 4. et 5., on considère une fonction $f$ deux fois dérivable sur l'intervalle
$[-4~;~2]$. On note $f'$ la fonction dérivée de $f$ et $f''$  la dérivée seconde de $f$.

On donne ci-dessous la courbe représentative $\mathcal{C}'$ de la fonction dérivée $f'$ dans un repère du plan. On donne de plus les points A$(-2~;~0)$, B(1~;~0) et C(0~;~5).

\begin{minipage}{0.52\linewidth}
\begin{enumerate}[resume]
\item La fonction $f$ est :


\begin{tabularx}{\linewidth}{*{2}{X}}
\textbf{a.~} concave sur $[-2~;~1]$ ; 	&\textbf{b.~} convexe sur $[-4~;~0]$;\\
\textbf{c.~}convexe sur $[-2~;~1]$;		&\textbf{d.~} convexe sur [0~;~2].
\end{tabularx}


\smallskip

\begin{tabular}{@{\hspace{0.05\linewidth}} | p{0.9\linewidth}}
La fonction $f'$ est croissante sur $[-4\;; 0]$ donc la fonction $f$ est convexe sur cet intervalle.
\hfill\textbf{Réponse b.}
\end{tabular}

\medskip

\item On admet que la droite (BC) est la tangente à la courbe $\mathcal{C}'$ au point B.
On a :


\begin{tabularx}{\linewidth}{*{2}{X}}
\textbf{a.~} $f'(1) < 0$ ;	&\textbf{b.~} $f'(1) = 5$ ;\\
\textbf{c.~} $f''(1) > 0$;	&\textbf{d.~} $f''(1) = - 5$.
\end{tabularx}


\begin{tabular}{@{\hspace{0.05\linewidth}} | p{0.9\linewidth}}
Le coefficient directeur de la droite (BC) est égal à $f''(1)$.
\hfill\textbf{Réponse d.}
\end{tabular}
\end{enumerate}
\end{minipage}\hfill
\begin{minipage}{0.45\linewidth}
\begin{center}
\psset{unit=1cm,arrowsize=2pt 3}
\begin{pspicture*}(-4.1,-3)(2,6)
\psaxes[linewidth=1.25pt,Dx=10,Dy=10]{->}(0,0)(-4,-3)(2,6)
\psaxes[linewidth=1.25pt](0,0)(1,1)
\psplot[plotpoints=2000,linewidth=1.25pt,linecolor=blue]{-4}{2}{x 2 add 1 x  sub mul 2.71828 x 0.5 mul exp mul}\uput[dr](-4,-1.5){\blue $\mathcal{C}'$}
\psline(-4,0.1)(-4,-0.1)
\uput[dr](-2,0){A} \uput[ur](1,0){B} \uput[l](0,5){C} 
\psplot[plotpoints=2000,linewidth=1.25pt]{-4}{2}{5 5 x  mul sub}
\uput[d](-4,0){$-4$}
\psdots(-2,0)(1,0)(0,5)
\end{pspicture*}
\end{center}

\end{minipage}

\begin{enumerate}[resume,start=6]
\item Soit $f$ la fonction définie sur $\R$ par $f(x) = \left(x^2 + 1\right)\text{e}^x$.

La primitive $F$ de $f$ sur $\R$  telle que $F(0) = 1$ est définie par:

\begin{center}
\begin{tabularx}{\linewidth}{*{2}{X}}
\textbf{a.~} $F(x) = \left(x^2 - 2x +3\right)\text{e}^x$ ;&\textbf{b.~} $F(x) = \left(x^2 - 2x + 3\right)\text{e}^x - 2;$\\
\textbf{c.~} $F(x) = \left(\dfrac13 x^3 + x\right)\text{e}^x + 1$ ;&\textbf{d.~} $F(x) = \left(\dfrac13 x^3 + x \right) \text{e}^x$.\\
\end{tabularx}
\end{center}

\begin{tabular}{@{\hspace{0.05\linewidth}} | p{0.9\linewidth}}
Pour la fonction $F$ de la réponse a. on a $F(0)=3$, et pour la fonction $F$ de la réponse d., on a $F(0)=0$. On peut donc éliminer ces deux réponses et tester les deux autres.

Si $F(x)= \left(x^2 - 2x + 3\right)\e^x - 2$, \\
alors $F'(x)=\left (2x-2\right )\e^{x} + \left(x^2 - 2x + 3\right)\e^x = \left(x^2 +1\right)\e^x  = f(x)$.

De plus $F(0)=3\e^{0}-2=1$.
\hfill\textbf{Réponse b.}
\end{tabular}

\end{enumerate}

\vspace{0.5cm}
